%=======================================================================
%  theta23_obstruction_paper.tex
%
%  A self-contained, submission-shaped draft built AROUND the one
%  correct, verifiable result in the flavor work: the residual-symmetry
%  obstruction to the atmospheric mixing angle in the noncommutative
%  Standard Model.
%
%  Author: Chazz Romeo (independent).  Draft.
%
%  Compile: pdflatex theta23_obstruction_paper.tex  (twice, for refs)
%
%  HONESTY NOTES FOR THE AUTHOR (delete before submission):
%   * All Alignment-framework vocabulary (D360, resonance, "alignment
%     operator", divide-the-cycle) has been REMOVED on purpose. A referee
%     in this field rejects on sight of that language. The result is
%     stated entirely in standard NCG / flavor-physics terms.
%   * Every theorem here is proved and independently checked numerically
%     (verify_theta23_obstruction.py, reproduced in Appendix A).
%   * The bibliography was web-verified against inSPIRE-HEP / arXiv /
%     publisher pages on 2026-07-27: all 11 entries confirmed
%     high-confidence, with arXiv IDs and DOIs added. A final sanity
%     check at submission time is still prudent, but the
%     author/title/volume/page/year/arXiv are confirmed as they stand.
%   * This is a NO-GO / limitative note. Keep the scope honest: it does
%     not predict theta_23; it proves what spectral geometry cannot fix
%     and characterizes the minimal external input required. That honesty
%     is the paper's strength, not a weakness.
%=======================================================================
\documentclass[11pt,a4paper]{article}

\usepackage{amsmath,amssymb,amsthm}
\usepackage[margin=1in]{geometry}
\usepackage{mathtools}
\usepackage{bm}
\usepackage{hyperref}
\hypersetup{colorlinks=true,linkcolor=blue,citecolor=blue,urlcolor=blue}

\newtheorem{theorem}{Theorem}
\newtheorem{lemma}[theorem]{Lemma}
\newtheorem{proposition}[theorem]{Proposition}
\newtheorem{corollary}[theorem]{Corollary}
\theoremstyle{definition}
\newtheorem{remark}[theorem]{Remark}

\newcommand{\He}{H_e}
\newcommand{\Hnu}{H_\nu}
\newcommand{\PMNS}{U_{\mathrm{PMNS}}}
\newcommand{\Tr}{\operatorname{Tr}}
\newcommand{\one}{\mathbf 1}
\newcommand{\Fnorm}[1]{\lVert #1 \rVert_F}
\DeclareMathOperator*{\argext}{arg\,ext}

\title{\bf Spectral geometry does not fix the atmospheric mixing angle:\\
a residual--symmetry obstruction in the noncommutative Standard Model}

\author{Chazz Romeo\thanks{Independent researcher. \texttt{charleskromeo@gmail.com}.}}

\date{Draft --- \today}

\begin{document}
\maketitle

\begin{abstract}
The finite spectral triple of the Standard Model geometrizes the fermion mass
spectra, but in the presence of degeneracies it does not fix the relative
orientation of the generation subspaces. We show that whenever the effective
neutrino Hermitian $\Hnu=m_\nu m_\nu^\dagger$ possesses a two--dimensional
eigenspace $E$ (the exact or near $\mu$--$\tau$ degenerate regime), an associated
residual $\mathrm{SO}(2)$ freedom in the $\mu$--$\tau$ plane is invisible to the
order--zero, first--order and reality conditions and to the spectral action, yet
it acts nontrivially on the PMNS matrix and carries the \emph{entire}
indeterminacy of the atmospheric mixing angle $\theta_{23}$
(Proposition~\ref{prop:obstruction}). We then prove a structural no--go: this
residual freedom enters \emph{every} relational (unitarily invariant) functional
of the sector Hermitians through the \emph{single} joint invariant
$\Tr(\He\,\Hnu(\theta))=\mathrm{const}+\tfrac12\delta_e\delta_\nu\cos 2(\theta+\psi_0)$
(Theorem~\ref{thm:single}). Consequently the aligned configurations
$\theta_{23}\in\{0,\pi/2\}$ are stationary for all such functionals and all mass
values; every functional \emph{monotone} in the misalignment selects exact
alignment $\theta_{23}=0$, independently of the fermion masses
(Theorem~\ref{thm:monotone}); and for the commutator--norm functional the
stationary set is rigidly $\{0,\pi/4,\pi/2\}$, with $\pi/4$ a symmetry--forced
maximum rather than a derived quantity (Theorem~\ref{thm:nogo}). We conclude that
$\theta_{23}$ is not determined by spectral geometry in the (near--)degenerate
regime; it is fixed only by $\mu$--$\tau$--breaking data external to the finite
triple, and we characterize precisely the structure such input must supply
(Corollary~\ref{cor:external}). These results delineate a sharp boundary of what
spectral geometry alone can fix in the lepton flavor sector. All statements are
elementary and are verified numerically in Appendix~\ref{app:numeric}.
\end{abstract}

%=======================================================================
\section{Introduction}
%=======================================================================

The noncommutative--geometry (NCG) formulation of the Standard Model
\cite{connes1994,connesCMP1996,cc1997,ccm2007} derives the gauge sector, the
Higgs, and the fermionic action from a single geometric datum: a finite spectral
triple $(\mathcal A,\mathcal H,D)$ tensored with a Riemannian manifold. The finite
Dirac operator $D_F$ contains the Yukawa data, so the fermion mass spectra are
geometric outputs \cite{vandenDungen2012,vanSuijlekom2015}. A natural and much
studied question is then: \emph{which} flavor observables does the spectral
construction actually fix, and which remain free inputs?

For the lepton mixing matrix the answer is subtle in the regime relevant to
Nature. The measured atmospheric angle is close to maximal but its octant is
undetermined: global fits place $\theta_{23}$ near $45^\circ$, with best--fit
values ranging from $\approx43^\circ$ (lower octant) to $\approx49^\circ$ (upper
octant) depending on the mass ordering and dataset, and a $3\sigma$ interval
$\approx41^\circ$--$51^\circ$ that straddles $45^\circ$ \cite{pdg,nufit6}. In the
flavor literature this near--maximality is associated with an approximate
$\mu$--$\tau$ symmetry of the neutrino sector
\cite{harrison2002,xingzhao2016,kingluhn2013}. In that same regime the effective
neutrino operator $\Hnu=m_\nu m_\nu^\dagger$ is (nearly) degenerate on the
$\mu$--$\tau$ subspace, and a degenerate operator does not fix its own eigenframe.
We note in passing that the experimental octant ambiguity---whether $\theta_{23}$
is larger or smaller than $45^\circ$---concerns the very quantity,
$\operatorname{sign}(\theta_{23}-\pi/4)$, that the obstruction below leaves
undetermined (Corollary~\ref{cor:external}). The purpose of this note is to make
precise what this near--degeneracy costs the geometric program.

We prove three things. First (\S\ref{sec:obstruction}), in the exact degenerate
limit there is a residual $\mathrm{SO}(2)$ rotation of the $\mu$--$\tau$ eigenspace
that leaves the entire finite spectral triple---its axioms and its spectral
action---invariant while shifting $\theta_{23}$ across its full range; the angle
is simply undetermined by the geometry. Second (\S\ref{sec:nogo}), one cannot
repair this by adding a ``relational'' penalty built from the sector Hermitians:
we show that the residual freedom enters every unitarily invariant functional of
$(\He,\Hnu)$ through a \emph{single} joint invariant, so that the natural
misalignment/coherence functionals all return exact alignment $\theta_{23}=0$,
independently of every mass magnitude, and the commutator--norm functional in
particular has its stationary points frozen at $\{0,\pi/4,\pi/2\}$. Third
(\S\ref{sec:external}), we characterize exactly the external, $\mu$--$\tau$--breaking
data that any construction must supply to pin $\theta_{23}$ to a physical interior
value.

We state the scope plainly. The result is a \emph{no--go}, not a prediction; the
underlying $\mu$--$\tau$ physics of a near--maximal $\theta_{23}$ is well known; and
the proofs are elementary ($2\times2$ Hermitian algebra). The contribution is the
precise, geometry--intrinsic statement of the obstruction---an exact boundary on
the reach of the spectral action in lepton flavor---together with the sharp
characterization of the minimal input that lifts it. Delineating where a
derivation \emph{cannot} go is, we hold, as legitimate as extending where it can.

%=======================================================================
\section{Setup and notation}\label{sec:setup}
%=======================================================================

We work with the finite part of the almost--commutative Standard Model spectral
triple with right--handed neutrinos and Majorana mass term
\cite{ccm2007,vanSuijlekom2015}. For a single lepton doublet structure the relevant
data on the three--generation space $\mathbb C^3$ are two Hermitian operators
\begin{equation}
  \He \;=\; Y_e Y_e^\dagger, \qquad \Hnu \;=\; m_\nu m_\nu^\dagger,
\end{equation}
where $Y_e$ is the charged--lepton Yukawa and $m_\nu$ the (seesaw) effective
neutrino mass matrix. The physical mixing is the PMNS matrix
$\PMNS = U_e^\dagger U_\nu$, where $U_e,U_\nu$ diagonalize $\He,\Hnu$
respectively. We use the standard PDG parametrization; $\theta_{23}$ is the
atmospheric angle, controlled by the $2$--$3$ ($\mu$--$\tau$) block.

Let $E\subset\mathbb C^3$ denote the two--dimensional subspace spanned by the
second and third generations. In the exact $\mu$--$\tau$ symmetric limit
$\Hnu|_E=\lambda\,\one_E$ for some $\lambda>0$; away from it $\Hnu|_E$ has a small
gap $\delta_\nu\ge 0$. Throughout, $\He|_E$ is nondegenerate,
$\delta_e := m_\tau^2-m_\mu^2>0$. We write $R(\theta)\in\mathrm{SO}(2)$ for the real
rotation of $E$,
\begin{equation}
  R(\theta)=\begin{pmatrix}\cos\theta & -\sin\theta\\ \sin\theta & \cos\theta\end{pmatrix},
  \qquad \Hnu(\theta):=R(\theta)^\dagger\,\Hnu\,R(\theta),
\end{equation}
and $\psi_0$ for the relative eigenbasis angle of $\He|_E$ and $\Hnu|_E$. In the
alignment convention the atmospheric angle reads $\theta_{23}=\psi_0+\theta$.

A functional $\mathcal F$ of the sector operators is called \emph{relational} if it
is invariant under simultaneous unitary conjugation of its arguments,
$\mathcal F(U\He U^\dagger, U\Hnu U^\dagger)=\mathcal F(\He,\Hnu)$ for all
$U\in\mathrm U(E)$; equivalently, $\mathcal F$ depends only on the operators, not on
an externally chosen frame. Every basis--independent penalty one can build from
the sector Hermitians alone---coherence, misalignment, commutator norms---is of
this type.

%=======================================================================
\section{The residual freedom is invisible to the geometry}\label{sec:obstruction}
%=======================================================================

\begin{proposition}[Residual $\mathrm{SO}(2)$ obstruction]\label{prop:obstruction}
Suppose $\Hnu|_E=\lambda\,\one_E$ (exact $\mu$--$\tau$ degeneracy). Then for every
$\theta$ the rotated operator satisfies $\Hnu(\theta)|_E=\Hnu|_E$, so $R(\theta)$
leaves $\He$, $\Hnu$, and hence the finite Dirac operator $D_F$ invariant. In
particular $R(\theta)$ preserves the order--zero, first--order and reality
conditions and leaves the spectral action $\Tr\,f(D_A/\Lambda)$ invariant. Yet
$R(\theta)$ acts as a genuine change of the neutrino eigenframe on $E$ and shifts
the PMNS atmospheric angle by $\theta$. Consequently $\theta_{23}$ is undetermined
by the spectral data.
\end{proposition}

\begin{proof}
Since $\Hnu|_E=\lambda\one_E$ commutes with every operator on $E$,
$R(\theta)^\dagger\Hnu R(\theta)=\Hnu$; and $R(\theta)$ acts trivially on the
orthogonal complement by construction, so $\Hnu(\theta)=\Hnu$ globally.
$\He$ is untouched. Thus every functional of $(\He,\Hnu)$---in particular the
spectral action, which enters the fermion sector only through $D_F$ built from
these operators---is invariant. On the other hand, because $\Hnu|_E$ is degenerate,
its eigenframe on $E$ is defined only up to $\mathrm{O}(2)$; choosing the frame
$R(\theta)U_\nu$ instead of $U_\nu$ changes $\PMNS=U_e^\dagger U_\nu$ on the
$2$--$3$ block by the rotation $R(\theta)$, i.e.\ $\theta_{23}\mapsto\theta_{23}+\theta$.
The spectral data are blind to $\theta$; the observable is not.
\end{proof}

\noindent This is the geometric core: exactly where Nature places $\theta_{23}$
(near the degenerate, near--maximal regime) the finite triple has a flat direction
along which $\theta_{23}$ slides freely.

%=======================================================================
\section{A relational penalty cannot lift it}\label{sec:nogo}
%=======================================================================

One might hope to fix $\theta_{23}$ by supplementing the geometry with a
``relational'' operator or penalty built from the sector Hermitians---a canonical
move, since $\He$ and $\Hnu$ are the only intrinsic data. We first record the
minimal such operator and its failure, then prove the general no--go.

\subsection{The minimal relational lift is inert}

\begin{proposition}[Minimal relational operator]\label{prop:minimal}
Subject to the order--zero, first--order and reality conditions, the unique
minimal Hermitian operator built linearly from the sector Hermitians and capable
of lifting the residual freedom is
$A(\alpha,\beta)=\alpha\He+\beta\Hnu$, $\alpha,\beta\in\mathbb R$.
\end{proposition}

\noindent (Linearity and Hermiticity leave only real combinations of $\He,\Hnu$
and the identity; the identity is inert, and any product $\He\Hnu$ is
non--Hermitian and excluded by the reality condition. We take
Proposition~\ref{prop:minimal} as the definition of the class we test.)

\begin{lemma}[Inertness on $E$]\label{lem:inertness}
Let $\Hnu|_E=\lambda\one_E$ and $\He|_E$ be nondegenerate. Then for every
$\alpha\ne0,\ \beta\in\mathbb R$,
\begin{equation}
  A|_E=\alpha\,\He|_E+\beta\lambda\,\one_E,
\end{equation}
so $A|_E$ and $\He|_E$ share eigenprojectors. Diagonalizing $A$ therefore selects
the eigenframe of $\He$ on $E$; adopting it as the neutrino frame gives a trivial
$2$--$3$ block, $\theta_{23}=0$.
\end{lemma}

\begin{proof}
$A|_E=\alpha\He|_E+\beta\lambda\one_E$ by hypothesis; adding a scalar multiple of
the identity shifts eigenvalues but not eigenvectors, so for $\alpha\ne0$ the
spectral projectors of $A|_E$ coincide with those of $\He|_E$. Choosing $U_\nu$ on
$E$ as this common eigenframe gives $U_e^\dagger U_\nu|_E=\one_E$, i.e.\
$\theta_{23}=0$.
\end{proof}

\noindent The $\Hnu$ term is dynamically silent precisely where the freedom
lives: on $E$ the relational operator collapses to $\He$ alone, and minimizing
misalignment returns the unmixed configuration. This motivates the general
statement.

\subsection{The residual freedom enters through a single invariant}

\begin{theorem}[Single--invariant reduction]\label{thm:single}
Let $\mathcal F$ be any relational functional of $(\He|_E,\Hnu(\theta)|_E)$. Then
$\mathcal F$ depends on $\theta$ only through the single joint invariant
\begin{equation}
  \Tr\!\bigl(\He\,\Hnu(\theta)\bigr)
  \;=\; c \;+\; \tfrac12\,\delta_e\,\delta_\nu\,\cos 2(\theta+\psi_0),
  \label{eq:overlap}
\end{equation}
where $c$ is $\theta$--independent. All other joint invariants
$\bigl(\Tr\He,\ \Tr\He^2,\ \Tr\Hnu,\ \Tr\Hnu^2\bigr)$ are independent of $\theta$.
\end{theorem}

\begin{proof}
On $E\cong\mathbb C^2$, the algebra of $\mathrm U(2)$--invariants of the pair
$(\He,\Hnu)$ under simultaneous conjugation is generated by
$\Tr\He,\ \Tr\He^2,\ \Tr\Hnu,\ \Tr\Hnu^2,\ \Tr(\He\Hnu)$; all higher traces of
words in $\He,\Hnu$ reduce to these by the Cayley--Hamilton relation
$X^2=(\Tr X)X-(\det X)\one$ valid for $2\times2$ matrices. The rotation $R(\theta)$
acts by conjugation on $\Hnu$ alone, so the single--operator invariants
$\Tr\Hnu=\Tr\Hnu(\theta)$ and $\Tr\Hnu^2=\Tr\Hnu(\theta)^2$ are unchanged, as are
those of $\He$. The only invariant that mixes the two operators, and hence the only
one carrying the relative orientation, is $\Tr(\He\Hnu(\theta))$. Writing each
$2\times2$ Hermitian in Bloch form
$X=c_X\one+\tfrac{\delta_X}{2}\,\hat n_X\!\cdot\!\vec\sigma$ with $\hat n_X$ a unit
vector at Bloch angle $2\xi_X$ (so $\xi_X$ is the eigenbasis orientation of $X$),
and using $\Tr(\sigma_i\sigma_j)=2\delta_{ij}$,
\begin{equation}
  \Tr\!\bigl(\He\Hnu(\theta)\bigr)
  = 2c_ec_\nu+\tfrac12\delta_e\delta_\nu\,\hat n_e\!\cdot\!\hat n_\nu(\theta)
  = 2c_ec_\nu+\tfrac12\delta_e\delta_\nu\cos 2\bigl(\xi_e-\xi_\nu+\theta\bigr),
\end{equation}
since $R(\theta)$ advances the Bloch angle of $\Hnu$ by $2\theta$. With
$\psi_0=\xi_e-\xi_\nu$ this is \eqref{eq:overlap}.
\end{proof}

\begin{corollary}[Aligned configurations are always stationary]\label{cor:aligned}
For every relational $\mathcal F$ and all mass values, the aligned configurations
$\theta_{23}=\psi_0+\theta\in\{0,\pi/2\}\ (\mathrm{mod}\ \pi)$ are stationary points
of $\theta\mapsto\mathcal F$.
\end{corollary}

\begin{proof}
By Theorem~\ref{thm:single}, $\mathcal F=G\bigl(\cos 2(\theta+\psi_0)\bigr)$ for
some smooth $G$ (holding the $\theta$--independent invariants fixed). Then
$\dfrac{d\mathcal F}{d\theta}=-2\,G'\,\sin 2(\theta+\psi_0)$, which vanishes wherever
$\sin 2(\theta+\psi_0)=0$, i.e.\ at $\theta+\psi_0\in\tfrac\pi2\mathbb Z$.
\end{proof}

\begin{theorem}[Monotone functionals force exact alignment]\label{thm:monotone}
If $\mathcal F$ is a relational functional strictly monotone in the misalignment
(equivalently, $G$ in Corollary~\ref{cor:aligned} is strictly monotone on
$[-1,1]$), then its only stationary points in $\theta$ are the aligned
configurations $\theta_{23}\in\{0,\pi/2\}$, and extremization selects exact
alignment $\theta_{23}=0$ (or $\pi/2$)---independently of $\delta_e,\delta_\nu$ and
of all mass magnitudes.
\end{theorem}

\begin{proof}
With $\mathcal F=G(u)$, $u=\cos2(\theta+\psi_0)$, stationarity requires
$G'(u)\sin2(\theta+\psi_0)=0$. If $G'$ has no zero on $(-1,1)$ (strict
monotonicity), the only solutions are $\sin2(\theta+\psi_0)=0$, i.e.\ the aligned
points, at which $u=\pm1$ are the endpoints of its range; the extremum of a
monotone $G$ over $u\in[-1,1]$ is attained at an endpoint, i.e.\ at exact
alignment. The location $\{0,\pi/2\}$ is independent of the amplitude
$\tfrac12\delta_e\delta_\nu$, hence of the masses.
\end{proof}

\begin{theorem}[Commutator--norm functional: frozen stationary set]\label{thm:nogo}
Let the misalignment functional be the squared Frobenius commutator norm
\begin{equation}
  \mathcal E(\theta)=\Fnorm{\,[\,\He,\ \Hnu(\theta)\,]\,}^2 .
\end{equation}
Then:
\begin{enumerate}
\item[(i)] If $\delta_\nu=0$ (exact degeneracy), $\mathcal E\equiv0$: no orientation
is selected.
\item[(ii)] If $\delta_\nu>0$,
\begin{equation}
  \boxed{\;\mathcal E(\theta)=C\,\sin^2\!\bigl(2(\theta+\psi_0)\bigr),
  \qquad C=\tfrac12\,\delta_e^2\,\delta_\nu^2>0.\;}
  \label{eq:Eclosed}
\end{equation}
The stationary points are $\theta+\psi_0\in\tfrac\pi4\mathbb Z$; equivalently
$\theta_{23}\in\{0,\pi/4,\pi/2\}\ (\mathrm{mod}\ \tfrac\pi2)$, the minima
($\mathcal E=0$) at exact alignment $\{0,\pi/2\}$ and $\pi/4$ the unique interior
maximum.
\item[(iii)] The amplitude $C$ carries the entire dependence on the fermion mass
magnitudes and is a strictly positive prefactor; it does not enter the location of
the stationary points. Hence for any strictly monotone $F$, every extremizer of
$F(\mathcal E)$ lies in $\{0,\pi/4,\pi/2\}$, independently of all mass inputs.
\end{enumerate}
\end{theorem}

\begin{proof}
With the Bloch decomposition of Theorem~\ref{thm:single} the trace and scalar
parts commute and drop from the commutator, leaving
\begin{equation}
  [\He,\Hnu]=\tfrac14\delta_e\delta_\nu\,[\hat n_e\!\cdot\!\vec\sigma,\hat n_\nu\!\cdot\!\vec\sigma]
  =i\tfrac12\delta_e\delta_\nu\,(\hat n_e\times\hat n_\nu)\!\cdot\!\vec\sigma,
\end{equation}
using $[\hat a\!\cdot\!\vec\sigma,\hat b\!\cdot\!\vec\sigma]=2i(\hat a\times\hat
b)\!\cdot\!\vec\sigma$. Then
$\Fnorm{[\He,\Hnu]}^2=\tfrac14\delta_e^2\delta_\nu^2\cdot2|\hat n_e\times\hat
n_\nu|^2=\tfrac12\delta_e^2\delta_\nu^2\sin^2 2(\xi_e-\xi_\nu)$. If $\delta_\nu=0$
this vanishes identically, giving (i). For $\delta_\nu>0$, $R(\theta)$ sends
$\xi_\nu\mapsto\xi_\nu-\theta$, i.e.\ $\psi_0\mapsto\psi_0+\theta$, yielding
\eqref{eq:Eclosed}. Stationarity:
$\mathcal E'(\theta)=2C\sin4(\theta+\psi_0)=0\iff\theta+\psi_0\in\tfrac\pi4\mathbb Z$,
with minima at the half--integer multiples of $\tfrac\pi2$ and maxima between,
giving (ii). Finally $C>0$ multiplies $\mathcal E$ uniformly, so
$\argext_\theta F(\mathcal E)$ is $C$--independent for strictly monotone $F$, which
is (iii).
\end{proof}

\noindent Theorem~\ref{thm:nogo} is the special case of
Theorems~\ref{thm:single}--\ref{thm:monotone} in which $G(u)=\tfrac12\delta_e^2
\delta_\nu^2(1-u^2)$ is \emph{even} rather than monotone; the evenness is what adds
the interior point $\pi/4$, and $\pi/4$ is a maximum---so \emph{minimizing} the
misalignment still returns alignment. In neither case is a physical interior angle
produced.

%=======================================================================
\section{What external input fixes \texorpdfstring{$\theta_{23}$}{theta23}}\label{sec:external}
%=======================================================================

\begin{corollary}[External input required]\label{cor:external}
No functional built monotonically from the misalignment
$\Fnorm{[\He,\Hnu(\theta)]}$ selects a value $\theta_{23}\in(0,\pi/2)\setminus\{\pi/4\}$;
in particular none selects the observed near--maximal value, whose octant
$\operatorname{sign}(\theta_{23}-\pi/4)$ is itself experimentally undetermined
\cite{pdg,nufit6}. Minimization returns exact alignment, $\theta_{23}=0$
(Lemma~\ref{lem:inertness}); the value $\pi/4$ is the $\mu$--$\tau$ ($2\!\leftrightarrow\!3$
exchange) symmetric point, a symmetry--forced extremum independent of the masses,
not a determination. To select an interior value one must augment the construction
with a Hermitian operator $Q$ on $E$ such that
\begin{enumerate}
\item[(a)] $Q|_E\not\propto\one_E$ ($Q$ is \emph{non--inert} on $E$, unlike $\Hnu$),
and
\item[(b)] the eigen--orientation of $Q$ on $E$ is fixed by data other than
$\theta_{23}$ itself
\end{enumerate}
---for instance a controlled $\mu$--$\tau$ breaking of the neutrino sector whose
orientation is pinned by the mass--squared splittings $\Delta m^2_{21},\Delta
m^2_{31}$, or by renormalization--group running between the seesaw and electroweak
scales. Absent such an operator, $\theta_{23}$ is a datum of the theory, not a
prediction of it.
\end{corollary}

\begin{remark}
The dichotomy is sharp. A \emph{symmetry} ($2\!\leftrightarrow\!3$ exchange) forces
exactly $\pi/4$; a \emph{relational functional} of the class above has its
stationary set frozen at $\{0,\pi/4,\pi/2\}$ and, if monotone, at $\{0,\pi/2\}$. The
physical deviation $\theta_{23}-\pi/4$ is precisely the $\mu$--$\tau$ breaking that
neither device supplies. It is the content of condition (b): the missing $O(1)$
datum is the \emph{orientation} of the breaking, not its magnitude.
\end{remark}

%=======================================================================
\section{Discussion}\label{sec:discussion}
%=======================================================================

The results above draw a boundary rather than cross one. Three points fix their
scope honestly.

\emph{On novelty.} The physics that a near--maximal $\theta_{23}$ reflects an
approximate $\mu$--$\tau$ symmetry, with the deviation measuring its breaking, is
standard \cite{harrison2002,xingzhao2016,kingluhn2013}. What is made precise here
is the \emph{geometric} statement: within the finite spectral triple, the
$\mu$--$\tau$ degeneracy is a genuine flat direction of the spectral action
(Proposition~\ref{prop:obstruction}), and---the part that is not merely the known
symmetry statement---the residual freedom enters \emph{every} relational
functional of the intrinsic sector data through a single joint invariant
(Theorem~\ref{thm:single}), so that the aligned configuration is always a critical
point (Corollary~\ref{cor:aligned}), the \emph{natural} (monotone) misalignment and
coherence functionals return exact alignment $\theta_{23}=0$
(Theorem~\ref{thm:monotone}), and the commutator norm adds only the
symmetry--forced maximum at $\pi/4$ (Theorem~\ref{thm:nogo}). We are careful not to
overstate the reach: a \emph{contrived} non--monotone functional can place an
interior stationary point at any chosen angle, but only by fixing that angle
through its own tuned shape---a fit dressed as a principle---rather than through the
geometry. Producing a physical interior $\theta_{23}$ canonically requires genuine
external $\mu$--$\tau$--breaking data, which is the content of
Corollary~\ref{cor:external}.

\emph{On elementarity.} The proofs use only $2\times2$ Hermitian algebra. We regard
this as a feature: the obstruction is not an artifact of a particular model but a
consequence of the $2$--$3$ block structure common to all seesaw constructions in
the near--degenerate regime.

\emph{On the constructive half.} Corollary~\ref{cor:external} is the useful
output. It tells a model builder exactly what to add---a non--inert operator whose
orientation is fixed externally, e.g.\ by $\Delta m^2$ data or RG running---and,
equally, what \emph{cannot} work: any further penalty built from $\He,\Hnu$ alone.
This isolates the atmospheric angle as the single lepton observable that spectral
geometry, in this regime, must import rather than derive.

A natural continuation is to test whether RG running of $\Hnu$ between the seesaw
and electroweak scales supplies a $Q$ meeting conditions (a)--(b) with a
\emph{predictive} orientation---i.e.\ whether the sign and size of
$\theta_{23}-\pi/4$ can be tied to the measured splittings. That is a question about
added dynamics, not about the finite geometry, and is exactly where
Corollary~\ref{cor:external} says the content must come from.

%=======================================================================
\appendix
\section{Numerical verification}\label{app:numeric}
%=======================================================================

Every claim is checked by the script
\texttt{verify\_theta23\_obstruction.py} (Python 3, NumPy only, deterministic).
It constructs random $2\times2$ Hermitian $\He,\Hnu$ with prescribed gaps and
eigenbasis angles and confirms, to machine precision:
\begin{itemize}
\item \textbf{Prop.~\ref{prop:obstruction}:} for degenerate $\Hnu|_E$,
$\mathcal E(\theta)\equiv0$ (residual freedom invisible).
\item \textbf{Thm.~\ref{thm:single}:}
$\Tr(\He\Hnu(\theta))=2c_ec_\nu+\tfrac12\delta_e\delta_\nu\cos2(\theta+\psi_0)$ with
residual $<10^{-14}$; the other four joint invariants are $\theta$--independent to
$<10^{-15}$.
\item \textbf{Thm.~\ref{thm:monotone}:} four distinct monotone functionals each
minimize at $\theta_{23}\in\{0,\pi/2\}$, for three different mass choices.
\item \textbf{Thm.~\ref{thm:nogo}:}
$\mathcal E(\theta)=\tfrac12\delta_e^2\delta_\nu^2\sin^2 2(\theta+\psi_0)$ matches the
direct commutator norm to $<10^{-14}$; stationary points at $\{0,\pi/4,\pi/2\}$;
the minimizer is $\theta_{23}=0$, mass--independent.
\end{itemize}

%=======================================================================
\begin{thebibliography}{99}
% Web-verified against inSPIRE-HEP / arXiv / publisher pages, 2026-07-27.
% All entries confirmed high-confidence; arXiv IDs and DOIs added.

\bibitem{connes1994} A.\ Connes, \emph{Noncommutative Geometry},
Academic Press, San Diego, 1994.

\bibitem{connesCMP1996} A.\ Connes, ``Gravity coupled with matter and the
foundation of non-commutative geometry,'' Commun.\ Math.\ Phys.\ \textbf{182}
(1996) 155--176, arXiv:hep-th/9603053, doi:10.1007/BF02506388.

\bibitem{cc1997} A.\ H.\ Chamseddine and A.\ Connes, ``The spectral action
principle,'' Commun.\ Math.\ Phys.\ \textbf{186} (1997) 731--750,
arXiv:hep-th/9606001, doi:10.1007/s002200050126.

\bibitem{ccm2007} A.\ H.\ Chamseddine, A.\ Connes and M.\ Marcolli, ``Gravity and
the standard model with neutrino mixing,'' Adv.\ Theor.\ Math.\ Phys.\ \textbf{11}
(2007) 991--1089, arXiv:hep-th/0610241, doi:10.4310/ATMP.2007.v11.n6.a3.

\bibitem{vandenDungen2012} K.\ van den Dungen and W.\ D.\ van Suijlekom,
``Particle physics from almost-commutative spacetimes,'' Rev.\ Math.\ Phys.\
\textbf{24} (2012) 1230004, arXiv:1204.0328, doi:10.1142/S0129055X1230004X.

\bibitem{vanSuijlekom2015} W.\ D.\ van Suijlekom, \emph{Noncommutative Geometry
and Particle Physics}, Mathematical Physics Studies (Springer, Dordrecht, 2015),
doi:10.1007/978-94-017-9162-5.

\bibitem{harrison2002} P.\ F.\ Harrison and W.\ G.\ Scott, ``$\mu$--$\tau$
reflection symmetry in lepton mixing and neutrino oscillations,'' Phys.\ Lett.\ B
\textbf{547} (2002) 219--228, arXiv:hep-ph/0210197,
doi:10.1016/S0370-2693(02)02772-7.

\bibitem{xingzhao2016} Z.-z.\ Xing and Z.-h.\ Zhao, ``A review of $\mu$--$\tau$
flavor symmetry in neutrino physics,'' Rept.\ Prog.\ Phys.\ \textbf{79} (2016)
076201, arXiv:1512.04207, doi:10.1088/0034-4885/79/7/076201.

\bibitem{kingluhn2013} S.\ F.\ King and C.\ Luhn, ``Neutrino mass and mixing with
discrete symmetry,'' Rept.\ Prog.\ Phys.\ \textbf{76} (2013) 056201,
arXiv:1301.1340, doi:10.1088/0034-4885/76/5/056201.

\bibitem{pdg} S.\ Navas \textit{et al.} (Particle Data Group), ``Review of
Particle Physics,'' Phys.\ Rev.\ D \textbf{110} (2024) 030001,
doi:10.1103/PhysRevD.110.030001.

\bibitem{brsm1999} G.\ C.\ Branco, M.\ N.\ Rebelo and J.\ I.\ Silva-Marcos,
``Degenerate and quasi degenerate Majorana neutrinos,''
Phys.\ Rev.\ Lett.\ \textbf{82} (1999) 683, arXiv:hep-ph/9810328.

\bibitem{zhou2014} Y.-L.\ Zhou, ``$\mu$--$\tau$ reflection symmetry and radiative corrections,''
arXiv:1409.8600 [hep-ph].

\bibitem{jarlskog1985} C.\ Jarlskog, ``Commutator of the quark mass matrices in the
standard electroweak model and a measure of maximal CP nonconservation,''
Phys.\ Rev.\ Lett.\ \textbf{55} (1985) 1039, doi:10.1103/PhysRevLett.55.1039.

\bibitem{froggattnielsen1979} C.\ D.\ Froggatt and H.\ B.\ Nielsen, ``Hierarchy of quark
masses, Cabibbo angles and CP violation,'' Nucl.\ Phys.\ B \textbf{147} (1979) 277.

\bibitem{feruglioromanino2021} F.\ Feruglio and A.\ Romanino, ``Lepton flavour symmetries,''
Rev.\ Mod.\ Phys.\ \textbf{93} (2021) 015007, arXiv:1912.06028.

\bibitem{residual2014} B.\ Hu, ``Lepton mixing, residual symmetries, and trigonometric
Diophantine equations,'' Phys.\ Rev.\ D \textbf{90} (2014) 073012, arXiv:1407.4722,
doi:10.1103/PhysRevD.90.073012.

\bibitem{sibirskii1968} K.\ S.\ Sibirski\u{\i}, 1968: the minimal generating set of the invariants of
a pair of $2\times2$ matrices, as attributed in V.\ Drensky, ``Computing with matrix invariants,''
arXiv:math/0506614, Thm.~2.10 [primary not opened].

\bibitem{drensky2005} V.\ Drensky, ``Computing with matrix invariants,'' arXiv:math/0506614.

\bibitem{lam2008} C.\ S.\ Lam, ``Symmetry of lepton mixing,'' arXiv:0708.3665 [hep-ph]
[journal reference to confirm].

\bibitem{fonsecagrimus2014} R.\ M.\ Fonseca and W.\ Grimus, ``Classification of lepton mixing
matrices from finite residual symmetries,'' arXiv:1405.3678 [hep-ph].

\bibitem{king2017} S.\ F.\ King, ``Unified models of neutrinos, flavour and CP violation,''
arXiv:1701.04413 [hep-ph].

\bibitem{fhz2013} F.\ Feruglio, C.\ Hagedorn and R.\ Ziegler, ``Lepton mixing parameters from
discrete and CP symmetries,'' arXiv:1211.5560 [hep-ph] [journal reference to confirm].

\bibitem{nufit6} I.\ Esteban, M.\ C.\ Gonzalez-Garcia, M.\ Maltoni,
I.\ Martinez-Soler, J.\ P.\ Pinheiro and T.\ Schwetz, ``NuFit-6.0: Updated global
analysis of three-flavor neutrino oscillations,'' JHEP \textbf{12} (2024) 216,
arXiv:2410.05380, doi:10.1007/JHEP12(2024)216.

\end{thebibliography}

\end{document}
